

\documentclass[notes, 11pt]{beamer} % THIS
% \documentclass[notes, 11pt]{beamer}
% \documentclass[notes, 11pt, aspectratio=169]{beamer}

% \usepackage{helvet}
% \usepackage[default]{lato}
% \usepackage{array}


% \documentclass[handout]{beamer}
% \usepackage{pgfpages}
% \pgfpagesuselayout{4 on 1}[a4paper,border shrink=5mm]

\usepackage{bbm}
\usepackage{amsthm}
\usepackage{amssymb}
\usepackage{xcolor}
\usepackage{graphicx}
\usepackage{epstopdf}
\usepackage{booktabs}
\usepackage[percent]{overpic}

\usepackage{sgamevar, tikz} % Game theory packages
\usetikzlibrary{trees, calc} % For extensive form games
\usepackage{subfig} % Manipulation and reference of small or sub figures and tables

\setbeamertemplate{theorems}[numbered] % to number
% \usetheme{Malmoe}
\usetheme{default}
% \usecolortheme{whale}
\usecolortheme{rose}


\newtheorem{defi}{Definition}
\newtheorem{teo}{Theorem}
\newtheorem{remi}{Remark}
\newtheorem{exi}{Example}
\newtheorem{propi}{Proposition}
\newtheorem{proteo}{Proof}
\newtheorem{lemi}{Lemma}
\newtheorem{prolemi}{Proof}
\newtheorem{propropi}{Proof}
\newtheorem{coro}{Corollary}
\newtheorem{proci}{Algorithm}
\newtheorem{ass}{Assumption}

\setbeamertemplate{footline}[frame number]

\usepackage{caption}
\captionsetup[figure]{labelformat=empty}%

\newenvironment{wideitemize}{\itemize\addtolength{\itemsep}{10pt}}{\enditemize}
% \usepackage{enumitem}

\usepackage[spanish]{babel}

% \renewcommand{\baselinestretch}{1}
% \expandafter\def\expandafter\insertshorttitle\expandafter{%
%   \insertshorttitle\hfill%
%   \insertframenumber\,/\,\inserttotalframenumber}

\DeclareMathOperator*{\argmax}{arg\,max}

\begin{document}
\title{Microeconom\'ia Aplicada (MAE)}   
\subtitle{Teor\'ia de Juegos: Juegos Estaticos con Informacion Completa}   
\author[Paz y Mino ]{Jos\'e Manuel Paz y Mi\~no}
\institute[shortinst]{FEN UCHILE}
\date{Oto\~no 2025} 

\frame{\titlepage} 



\begin{frame}{Dragu, T. (2011)}
	\begin{wideitemize}
		\item<1-> Motivaci\'on: 
		\begin{wideitemize}
		\item<2->[--] Ataques terroristas incrementan vigilancia por parte de gobiernos.
		\item<3->[--] Esta vigilancia se traduce en menor privacidad y mayor monitoreo de actividades.
		\item<4->[--] Menos privacidad genera menos terrorismo?
		\end{wideitemize}
		\item<5-> Resultados
		\begin{wideitemize}
			\item<6->[--] Menor privacidad puede incrementar actividad terrorista.
			\item<7->[--] Una agencia antiterrorista siempre prefiere menos privacidad.
		\end{wideitemize}
	\end{wideitemize}
\end{frame}

\begin{frame}{Dragu, T. (2011)}
	\begin{wideitemize}
		\item<1-> Mecanismos:
		\begin{wideitemize}
			\item<2-> Efectos directos (costos) e indirectos (estrat\'egicos) de mayor privacidad.
			\item<3-> Para terroristas ambos efectos van en misma direcci\'on.
			\item<4-> Para agencia antiterrorista van en sentidos opuestos.
		\end{wideitemize}

	\end{wideitemize}
\end{frame}





\begin{frame}
\maketitle
\end{frame}

\end{document}
